
\documentclass[11pt]{letter}

\usepackage[bitstream-charter]{mathdesign}
\usepackage{microtype, graphicx}
\usepackage[letterpaper, margin=1in]{geometry}

\longindentation=30pt

\signature{Johanna H\"aussler, Gy\"orgy Barab\'as \& Anna Ekl\"of}

\address{Johanna H\"aussler \\ German Centre for Integrative Biodiversity Research (iDiv) \\ Halle-Jena-Leipzig \\ Deutscher Platz 5e, 04103 Leipzig, Germany}

\date{January 17, 2020}


\begin{document}

\begin{letter}{
    \vspace{-2cm}
    
    Dr.~Tim Coulson\\
    Editor-in-Chief\\
    Ecology Letters}


\opening{Dear Dr.~Coulson,}

Please find enclosed our manuscript ``A  Bayesian network approach to trophic metacommunities shows habitat loss accelerates top species extinctions'', which we would like to be considered for publication in Ecology Letters as a Letter.

Habitat loss is one of the main drivers causing species extinction rates to accelerate, yet its effect on complex communities such as food webs remains mostly unknown. We address this issue by developing a novel approach to studying trophic metacommunities. Our method combines classic metapopulation models on fragmented landscapes with a Bayesian network representation of trophic interactions for calculating local extinction rates. Our approach offers several advantages over existing methods. First, due to its computational efficiency, it can be applied to much larger food webs and landscapes than would be feasible with fully-fledged dynamical models. Second, one can obtain analytical solutions for simple but important food web structures, such as a linear food chain. Third, since the method is built on classical metapopulation theory, many of the theory's known properties can be repurposed to aid the analysis of trophic metacommunities---such as being able to rank the habitat patches of the landscape with respect to their importance to the persistence of the metacommunity as a whole.

Here we make use of this patch ranking to study how progressive habitat loss affects species extinctions depending on whether one prioritizes the removal of valuable vs.\ non-valuable patches for the metacommunity. The following key results emerge:

\begin{itemize}

\item Species at higher trophic levels suffer elevated extinction risks, to a degree that this itself puts a strict limit on the maximum number of trophic levels that can persist, irrespective of energetic or other constraints.

\item Trophic metacommunities can tolerate substantial habitat loss without extinctions if the least valuable patches are removed first. By contrast, starting removal with the most important patches greatly accelerates metacommunity collapse. This underlines the necessity of estimating patch values in real-life conservation efforts---which our approach provides a method for.

\item Surprisingly, and importantly from a conservation perspective, removing patches in random order is almost indistinguishable in its effects from the worst-case scenario of removing the most valuable patches first.

\end{itemize}

Exploring the effect of habitat loss on complex communities is an important and timely topic in community ecology. Given the general interest of the question we address, and trying to reach as wide a readership as possible, we believe Ecology Letters would be an ideal venue for our manuscript.
% Given the novelty of our method and its broad implications, we create an openly available paradigm for studying trophic metacommunities, fitting well with the journal's profile. Our method combines the flexibility of explicit dynamical models with the tractability and computational efficiency to simple topological methods, thereby providing an alternative, complementary way of analyzing spatial food webs. 

We thank you for your dedicated attention to our work.

This manuscript has not been published and is not under consideration for publication elsewhere. We have no conflicts of interest to disclose.

\closing{Sincerely,}

\end{letter}
\end{document}

